\subsection*{NT-A15--A17. Several Hecke coordinates and the exact independence gate}

Retain \(L/K,G,C,g,H,F,h,\delta,n_L,S,\mathfrak s_F\)
from NT-A12. Fix \(r\ge1\) unitary Hecke characters
\(\eta_1,\ldots,\eta_r\) of \(K\). Require that the character
\[
             \prod_{i=1}^r\eta_i^{m_i},
                    \qquad 0\ne\mathbf m\in\mathbf Z^r,  \tag{A15.1}
\]
is never finite order times a unitary norm power.
Enlarge \(S\) to contain all their conductor primes.
Write \(\nu_i=\eta_i\circ N_{F/K}\), and put
\[
 \begin{aligned}
 \mathfrak f_*&=\operatorname{lcm}_i\mathfrak f_{\nu_i},\\
 B_*&=4+\sum_{i=1}^r\max_j|\kappa_{\nu_i,j}|,\\
 Q_*&=D_L(N_F\mathfrak f_*\,N_F\mathfrak s_F)^h B_*^{n_L},\\
 {\cal L}_*(x)&=\log Q_*+2n_L\log(3x).
 \end{aligned}                                             \tag{A15.2}
\]
The exact RH family is
\[
 L_F\left(s,\chi\prod_{i=1}^r\nu_i^{m_i}\right),
       \qquad \chi\in\widehat H,\quad\mathbf m\in\mathbf Z^r .
                                                               \tag{A15.3}
\]
It is a subfamily of \(M\). The mixed products in this family
are part of the hypothesis; separate assertions only for
the individual coordinates are not substituted for them.

\paragraph{Classical discrepancy input, with a precise locator.}
For \(N\) points \(\boldsymbol\theta_j\in[0,1)^r\), the
Erd\H{o}s--Tur\'an--Koksma inequality bounds their normalized
rectangular discrepancy by
\[
 \begin{gathered}
 R(\mathbf m)=\prod_i\max(1,|m_i|),\\
 D_N\le(3/2)^r\left\{\frac2{H+1}
     +\sum_{0<\|\mathbf m\|_\infty\le H}
       \frac{1}{R(\mathbf m)}
       \left|\frac1N\sum_{j=1}^N
                  e^{2\pi i\mathbf m\cdot\boldsymbol\theta_j}
       \right|\right\}.
 \end{gathered}                                             \tag{A15.4}
\]
The exact statement is Dick--Pillichshammer,
\emph{Digital Nets and Sequences}, author preprint
Theorem 3.26, printed/PDF p.~84 (S-NT-DP).
That page refers its proof to Drmota--Tichy,
\emph{Sequences, Discrepancies and Applications},
\S1.2.2. The theorem page was image checked; the
referred proof is a cited premise, not a newly audited
argument in this run.

\paragraph{NT-A15: joint boxes and a least prime.}
For a product \(J=J_1\times\cdots\times J_r\) of circular
arcs, let \(\mu=\prod_i|J_i|\). Under (A15.1)--(A15.3),
the count of primes outside \(S\), Frobenius class \(C\),
norm at most \(x\), and coordinate angles in \(J\) is
\[
 N_{C,J}^{(r)}(x)=\delta\mu\operatorname{Li}_2(x)
       +O_r\bigl(\delta\sqrt x
                    (\log(2x))^r{\cal L}_*(x)\bigr).
                                                               \tag{A15.5}
\]
The constant depends only on \(r\); all other dependence
is displayed. Endpoint conventions are arbitrary, and
the same bound holds for absolute-degree-one primes.
For \(\mu>0\), this box contains such a degree-one prime
of norm at most
\[
 C_r\mu^{-2}V_*^{\,2r+4},\qquad
       V_*=\log(eQ_*)+n_L\log(e n_L/\mu),                 \tag{A15.6}
\]
where \(C_r\) depends only on \(r\).

\begin{proof}
Apply the diagonal-idele argument of NT-A12 to each
nonzero monomial in (A15.1). Its norm pullback is still
not finite order times norm. Multiplication by a
finite-order \(\chi\) therefore cannot make it pure norm.
Only \(\chi=1,\mathbf m=0\) is principal.
For negative powers the finite conductor is unchanged;
for any product its conductor divides the least common
multiple of the individual conductors, since all factors
are trivial on each sufficiently deep local unit group.
Thus
\[
 {\cal C}\left(\chi\prod_i\nu_i^{m_i}\right)
  \le D_FN_F\mathfrak f_\chi N_F\mathfrak f_*
       [B_*(1+\|\mathbf m\|_\infty)]^{n_F}.              \tag{A15.7}
\]
At a complex place the absolute value of the sum of
angular parameters is at most the sum of their absolute
values; imaginary parameters add. At real places the
parity remains zero or one. These facts prove the
gamma-factor part of (A15.7), including negative \(m_i\).
Summing logarithms over \(\chi\) uses exactly
\(D_L=D_F^h\prod_\chi N_F\mathfrak f_\chi\).

The whole cyclic induction and omitted-prime argument
(A12.8)--(A12.10) now applies to the monomial.
In particular its relative-degree correction remains
\(\delta n_L\sqrt x\), independently of \(\mathbf m\).
It gives the complete Fourier estimate
\[
 \begin{split}
 T_{\mathbf m}^C(x)
  &:=
    \sum_{\substack{N_K\mathfrak p\le x,\ \mathfrak p\notin S\\
                     \operatorname{Frob}_{\mathfrak p}\in C}}
                          \prod_i\eta_i(\mathfrak p)^{m_i}\\
  &=\delta\mathbf1_{\mathbf m=0}\operatorname{Li}_2(x)
      +O\left(\delta\sqrt x\{
         \log Q_*+n_L\log[(1+\|\mathbf m\|_\infty)(2x)]
                                  \}\right).
 \end{split}                                                   \tag{A15.8}
\]
The constant here is absolute: any dependence on \(r\)
has been put in \(B_*\), not in the Hecke prime estimate.
The principal factor is again handled separately.

Use (A15.4) in unnormalized form with \(H=\lceil x\rceil\).
The sum of the frequency weights is exactly
\[
 \sum_{0<\|\mathbf m\|_\infty\le H}\frac1{R(\mathbf m)}
       =\left(1+2\sum_{j=1}^H\frac1j\right)^r-1
                     \ll_r(\log(2H))^r .
 \]
The zero-frequency truncation contributes
\(O_r([K:\mathbf Q])\), absorbed as in NT-A12 by
\(\delta n_L\le\delta\log Q_*/\log4\).
All the braces in (A15.8) are \(O({\cal L}_*(x))\)
on this frequency cube. These facts prove (A15.5)
for half-open boxes not crossing the circle origins.
A circular box is a disjoint union of at most \(2^r\)
such boxes. Shrinking a coordinate interval onto a
specified endpoint bounds each coordinate-face count
by the same error. Summing the finitely many faces
permits all endpoint conventions.
The number of omitted absolute-degree-at-least-two
primes is at most \([K:\mathbf Q]\sqrt x\), again
absorbed. If the prime list is empty, the unnormalized
discrepancy assertion follows directly from the
zero-frequency estimate, so no division by zero is used.

For (A15.6), set \(V=V_*\), \(n=n_L\), and
\(x=C_r\mu^{-2}V^{2r+4}\). The inequalities
\(V\ge n\), \(n\log(1/\mu)\le V\), and
\(n\log V\le2V\), proved in NT-A14, imply
\[
 \begin{aligned}
 \log x&\le(\log C_r+4r+10)V/n,\\
 \log(2x)&\le(\log C_r+4r+11)V/n,\\
 {\cal L}_*(x)&\le(8r+24+2\log C_r)V.
 \end{aligned}
\]
Let \(A_r\) be a constant in the degree-one error
(A15.5). Since \(\operatorname{Li}_2(x)\ge x/(2\log x)\),
positivity follows, after cancellation of \(\delta\),
from
\[
 \sqrt{C_r}>
  2A_r(\log C_r+4r+10)(\log C_r+4r+11)^r
                    (8r+24+2\log C_r).
 \]
An appropriate \(C_r\ge4\) exists for every fixed \(r\).
The discarded factor \(n^{-(r+1)}\) is at most one.
This proves the least-prime assertion.
\end{proof}

\paragraph{NT-A16: simultaneously shrinking boxes.}
Fix \(c_i>0\), \(\alpha_i>0\), and put
\[
 A=\sum_i\alpha_i<\tfrac12,\qquad
 c_*=\prod_i c_i,\qquad
 u_0=\max\{2,c_1^{1/\alpha_1},\ldots,c_r^{1/\alpha_r}\}.
 \]
Fix every coordinate center and count \(u_0<N_K\mathfrak p\le x\)
in Frobenius class \(C\) whose \(i\)-th angular coordinate
lies in the arc of length \(c_i(N_K\mathfrak p)^{-\alpha_i}\).
Then
\[
 \begin{split}
 N_{C,\mathbf c,\boldsymbol\alpha}(x)
   ={}&\delta c_*\int_{u_0}^x\frac{u^{-A}}{\log u}\,du\\
     &+O_{\mathbf c,\boldsymbol\alpha,r}\left(
        \delta\sqrt x(\log(2x))^r{\cal L}_*(x)\right.\\
     &\hspace{38mm}\left.
        +\delta x^{3/4-A/2}
             (\log(2x))^{(r-1)/2}\sqrt{{\cal L}_*(x)}
                                    \right).
 \end{split}                                                   \tag{A16.1}
\]
The error is uniform in the centers and all arithmetic
parameters already displayed. The degree-one version
holds too. A sufficient relative asymptotic regime is
\[
          {\cal L}_*(x)(\log x)^{r+1}
                                  =o(x^{1/2-A}).           \tag{A16.2}
\]

\begin{proof}
The boxes are nested because every side length decreases.
Partition \((u_0,x]\) with multiplicative ratio
\(1+\varepsilon\), \(0<\varepsilon\le1/2\).
Applying (A15.5) to both endpoints gives a total error
at most a constant times \(A_0/\varepsilon\), where
\[
 A_0=\delta\sqrt x(\log(2x))^r{\cal L}_*(x).
 \]
Indeed the endpoint square roots form the same geometric
series as in NT-A13. The volume is exactly \(c_*u^{-A}\);
its ratio on a mesh interval is at most
\((1+\varepsilon)^A=1+O_A(\varepsilon)\).
Consequently the integral bracketing error is
\(O_{\mathbf c,\boldsymbol\alpha}(B_0\varepsilon)\), with
\[
                  B_0=\frac{\delta c_* x^{1-A}}{\log x}.
 \]
The canonical \(u_0\) bounds the finite initial range
using only the fixed \(\mathbf c,\boldsymbol\alpha\);
as before its field degree is absorbed by
\([K:\mathbf Q]\le\delta n_L\).
Choose \(\varepsilon=\min\{1/2,\sqrt{A_0/B_0}\}\).
The resulting error is \(O(A_0+\sqrt{A_0B_0})\), and
\[
 \sqrt{A_0B_0}
  =\delta\sqrt{c_*}\,x^{3/4-A/2}
          \sqrt{\frac{(\log(2x))^r}{\log x}{\cal L}_*(x)}.
 \]
This is the second term of (A16.1). Its main term is
asymptotic to
\(\delta c_*x^{1-A}/[(1-A)\log x]\).
The two relative errors are bounded by fixed constants
times \(R+\sqrt R\), where
\(R={\cal L}_*(x)(\log x)^{r+1}x^{A-1/2}\).
This proves (A16.2), and the degree-one argument is
identical. No assertion is made at \(A=1/2\).
\end{proof}

\paragraph{NT-A17: necessity of the monomial independence condition.}
For fixed \(L/K,C,S\), under \(M\) the joint prime angles
in this conjugacy class are Haar equidistributed on
\((\mathbf R/\mathbf Z)^r\) if and only if (A15.1) holds.
More precisely, if
\(\prod_i\eta_i^{m_i}=\theta|\cdot|^{it}\) for nonzero
\(\mathbf m\), with \(\theta\) of finite order dividing \(b\),
rename the real norm parameter so that the ideal value of
the \(b\)-th power is \(q^{i\tau}\). Then its normalized
prime Fourier moment is
\[
 \frac{1}{N_C(x)}
 \sum_{\substack{N_K\mathfrak p\le x,\ \mathfrak p\notin S\\
                  \operatorname{Frob}_{\mathfrak p}\in C}}
          \prod_i\eta_i(\mathfrak p)^{b m_i}
       =\frac{x^{i\tau}}{1+i\tau}+o(1).                  \tag{A17.1}
\]
Here \(N_C(x)\) is the number of primes in the sum.
For \(\tau=0\) the left side is exactly one once \(N_C(x)>0\).
Thus neither separate infinite order nor separate
coordinate equidistribution replaces (A15.1).

\begin{proof}
Sufficiency follows from (A15.5) for every fixed box:
its error is \(o(\delta\operatorname{Li}_2(x))\) for fixed
data, and the zero-frequency count has that denominator
asymptotic. Boxes characterize Haar equidistribution.

For necessity, raise the proposed relation to the \(b\)-th
power. Outside the conductor set, every summand in (A17.1)
is exactly \((N_K\mathfrak p)^{i\tau}\).
The fixed-field Chebotarev estimate under RH for \(\zeta_L\),
which \(M\) supplies, gives
\[
 N_C(u)=\delta\operatorname{Li}_2(u)
                       +O_{L,K,C,S}(\sqrt u\log(2u)).
 \]
Abel summation, with fixed \(\tau\), therefore yields
\[
 \sum_{\mathfrak p\text{ in the sum}}
                  (N_K\mathfrak p)^{i\tau}
    =\delta\int_2^x\frac{u^{i\tau}}{\log u}\,du
                +O_{L,K,C,S,\tau}(\sqrt x\log(2x)).
 \]
The lower endpoint can add only a fixed constant.
Differentiating
\(u^{1+i\tau}/[(1+i\tau)\log u]\) shows that this integral is
\[
       \frac{x^{1+i\tau}}{(1+i\tau)\log x}
                         +O_\tau(x/(\log x)^2).
 \]
Division by \(N_C(x)\sim\delta x/\log x\) proves (A17.1).
Its modulus tends to \(1/\sqrt{1+\tau^2}>0\).
The corresponding nonzero Fourier character has Haar
integral zero, so Haar equidistribution fails.
\end{proof}

\paragraph{Geometric meaning and work placement.}
The character lattice \(\mathbf Z^r\) describes the
compact torus coordinates; (A15.1) removes every
finite character or norm drift in that lattice.
For normalized algebraic Hecke characters this is the
torus part of the tensor category of the corresponding
CM objects, coupled to the finite \'{e}tale \(G\)-cover.
The mixed Hecke family (A15.3) records every tensor
character needed by the proof. This is not statistical
independence of zeros or a new assertion of functoriality
for general automorphic representations.
Full proofs: \texttt{joint\_hecke\_tori.tex}, NT-A15--A17,
also \texttt{section.tex}; run
\texttt{GRH\_NIGHT/number\_theory}.
Classical premises are NT-A08/NT-A12, S-NT-NEUKIRCH,
S-NT-DP and the explicit Chebotarev result NT-C02.
